%! Sine and Cosine Function Values — Unit 6, Lesson 6.4 — Group Activity (Answer Key)
\documentclass[10pt]{article}
\usepackage{precalculus-article}
\usepackage{precalculus-key}   % pulls in -boxes; do NOT also load -boxes
\newcommand{\TallMath}[1]{$\displaystyle #1\rule[-1.4em]{0pt}{3.2em}$}

\begin{document}

\pageheader{Unit 6, Lesson 6.4}{Group Activity --- Answer Key}
\namepartnerperiod

\vspace{0.1in}

\begin{scenariobox}[Exact Values in Engineering and Navigation]{sky}
\small
Engineers and navigators routinely need exact trigonometric values rather than decimal approximations.
In this activity your group will build fluency with exact values and apply them to a real-world context.
\end{scenariobox}

\vspace{0.1in}

%-----------------------------------------------------------------------
% TIER R
%-----------------------------------------------------------------------
\begin{tcolorbox}[colback=white, colframe=black!40,
    title=\textbf{Tier R --- Build the Unit Circle Table},
    fonttitle=\bfseries, arc=2mm, left=3mm, right=3mm, top=2mm, bottom=2mm]
\small
\renewcommand{\arraystretch}{2.0}
\begin{tabularx}{\linewidth}{@{} >{\centering\arraybackslash}p{2.2cm}
    >{\centering\arraybackslash}X >{\centering\arraybackslash}X @{}}
\toprule
\textbf{Angle $\theta$} & \textbf{$\cos\theta$} & \textbf{$\sin\theta$} \\
\midrule
$0$      & $1$              & $0$              \\
$\pi/6$  & $\sqrt{3}/2$     & $1/2$            \\
$\pi/4$  & $\sqrt{2}/2$     & $\sqrt{2}/2$     \\
$\pi/3$  & $1/2$            & $\sqrt{3}/2$     \\
$\pi/2$  & $0$              & $1$              \\
$2\pi/3$ & \ans{$-1/2$}             & \ans{$\sqrt{3}/2$}       \\
$3\pi/4$ & \ans{$-\sqrt{2}/2$}      & \ans{$\sqrt{2}/2$}       \\
$5\pi/6$ & \ans{$-\sqrt{3}/2$}      & \ans{$1/2$}              \\
$\pi$    & \ans{$-1$}               & \ans{$0$}                \\
\bottomrule
\end{tabularx}
\renewcommand{\arraystretch}{1}

\vspace{0.08in}
\textbf{Reflection:} \ans{$\sin\theta$ increases from 0 to 1 on $[0, \pi/2]$, then decreases back to 0 on $[\pi/2, \pi]$ --- it is symmetric about $\theta = \pi/2$.}
\end{tcolorbox}

\vspace{0.08in}

%-----------------------------------------------------------------------
% TIER A
%-----------------------------------------------------------------------
\begin{tcolorbox}[colback=white, colframe=black!40,
    title=\textbf{Tier A --- Ferris Wheel Coordinates},
    fonttitle=\bfseries, arc=2mm, left=3mm, right=3mm, top=2mm, bottom=2mm]
\small
$r = 40$~ft. Apply $x = 40\cos\theta$, $y = 40\sin\theta$.

\renewcommand{\arraystretch}{2.0}
\begin{tabularx}{\linewidth}{@{} >{\centering\arraybackslash}p{2.0cm}
    >{\centering\arraybackslash}X >{\centering\arraybackslash}X @{}}
\toprule
\textbf{Angle $\theta$} & \textbf{$x$ (ft)} & \textbf{$y$ (ft)} \\
\midrule
$\pi/6$  & \ans{$20\sqrt{3}$} & \ans{$20$}           \\
$\pi/4$  & \ans{$20\sqrt{2}$} & \ans{$20\sqrt{2}$}   \\
$\pi/3$  & \ans{$20$}         & \ans{$20\sqrt{3}$}   \\
$\pi/2$  & \ans{$0$}          & \ans{$40$}           \\
$2\pi/3$ & \ans{$-20$}        & \ans{$20\sqrt{3}$}   \\
$5\pi/6$ & \ans{$-20\sqrt{3}$}& \ans{$20$}           \\
\bottomrule
\end{tabularx}
\renewcommand{\arraystretch}{1}

\vspace{0.08in}
\textbf{Question:} \ans{The car is highest at $\theta = \pi/2$, where $y = 40$~ft. The maximum possible height equals the radius, 40~ft (top of the wheel).}
\end{tcolorbox}

\vspace{0.08in}

%-----------------------------------------------------------------------
% TIER E
%-----------------------------------------------------------------------
\begin{tcolorbox}[colback=white, colframe=black!40,
    title=\textbf{Tier E --- Symmetry and Inverse Reasoning},
    fonttitle=\bfseries, arc=2mm, left=3mm, right=3mm, top=2mm, bottom=2mm]
\small
$P = \left(-\dfrac{r\sqrt{3}}{2},\, \dfrac{r}{2}\right)$

\textbf{(a)} $\cos\theta = -\tfrac{\sqrt{3}}{2}$ and $\sin\theta = \tfrac{1}{2}$.
$\sin\theta > 0$ places us in QI or QII; $\cos\theta < 0$ places us in QII or QIII. Therefore $\theta$ is in \textbf{QII}.
Reference angle: $\cos^{-1}\!\left(\tfrac{\sqrt{3}}{2}\right) = \pi/6$, so $\theta = \pi - \pi/6$.

$\theta = $ \ans{$\dfrac{5\pi}{6}$}

\vspace{0.08in}
\textbf{(b)} On the unit circle, $\sin\alpha = \tfrac{1}{2}$. This occurs at $\alpha = \pi/6$ (QI) and $\alpha = 5\pi/6$ (QII), since the unit circle is symmetric about the $y$-axis.

$\alpha = $ \ans{$\dfrac{\pi}{6}$} and $\alpha = $ \ans{$\dfrac{5\pi}{6}$}

\vspace{0.06in}
\textbf{Symmetry explanation:}

\begin{teachernote}
Students should articulate that the horizontal line $y = \tfrac{1}{2}$ intersects the unit circle at exactly two points, one in QI and one in QII, related by reflection across the $y$-axis. This gives the identity $\sin(\pi - \alpha) = \sin\alpha$. Credit answers that mention ``reflection,'' ``symmetry about the $y$-axis,'' or the identity explicitly.
\end{teachernote}

\ans{The line $y = \tfrac{1}{2}$ intersects the unit circle in two points. They are reflections of each other across the $y$-axis, so if one is at angle $\alpha$ then the other is at $\pi - \alpha$. This is why $\sin(\pi/6) = \sin(5\pi/6) = \tfrac{1}{2}$.}
\end{tcolorbox}

\end{document}
